\chapter{Auxiliary mathematical definitions}
\label{app:aux-defs}

\section*{Overview}

This appendix collects the mathematical vocabulary that the main
text uses without re-deriving --- inner products, completeness,
dual spaces, adjoints, tensor products. Readers comfortable with
elementary linear algebra and functional analysis can skip it;
others may want to skim it once before \cref{ch:postulates} or
refer back to it as the chapters cite it.

% --------------------------------------------------------------
\section{Forms, inner products, and norms}
\label{app:inner-product}

\begin{definition}[Sesquilinear form]
\label{def:sesquilinear}
Let $V$ be a complex vector space. A map
$\langle\cdot,\cdot\rangle:V\times V\to\C$ is \emph{sesquilinear}
if it is linear in one argument and conjugate-linear (antilinear)
in the other. We follow the physics convention of antilinearity in
the first slot and linearity in the second: for all
$\alpha,\beta\in\C$ and $\phi,\psi,\chi\in V$,
\begin{equation}
\langle\alpha\phi+\beta\psi,\chi\rangle
= \bar\alpha\langle\phi,\chi\rangle + \bar\beta\langle\psi,\chi\rangle,
\qquad
\langle\phi,\alpha\psi+\beta\chi\rangle
= \alpha\langle\phi,\psi\rangle + \beta\langle\phi,\chi\rangle.
\end{equation}
\end{definition}

\begin{definition}[Conjugate symmetry]
\label{def:conj-sym}
A form $\langle\cdot,\cdot\rangle:V\times V\to\C$ is
\emph{conjugate-symmetric} if
$\langle\phi,\psi\rangle=\overline{\langle\psi,\phi\rangle}$ for
all $\phi,\psi\in V$. In particular
$\langle\psi,\psi\rangle\in\R$.
\end{definition}

\begin{definition}[Positive definiteness]
\label{def:positive-definite}
A form $\langle\cdot,\cdot\rangle$ on $V$ is
\emph{positive definite} if $\langle\psi,\psi\rangle\geq0$ for
every $\psi\in V$, with equality if and only if $\psi=0$.
\end{definition}

\begin{definition}[Inner product, norm, metric]
\label{def:inner-product}
An \emph{inner product} on a complex vector space $V$ is a
sesquilinear (\cref{def:sesquilinear}), conjugate-symmetric
(\cref{def:conj-sym}), positive-definite
(\cref{def:positive-definite}) form
$\langle\cdot,\cdot\rangle:V\times V\to\C$. It induces a
\emph{norm} $\|\psi\|=\sqrt{\langle\psi,\psi\rangle}$ and, through
the norm, a \emph{metric} $d(\phi,\psi)=\|\phi-\psi\|$.
\end{definition}

% --------------------------------------------------------------
\section{Completeness and separability}
\label{app:completeness}

\begin{definition}[Cauchy sequence and completeness]
\label{def:cauchy}
A sequence $(\psi_n)$ in a normed space $V$ is a \emph{Cauchy
sequence} if for every $\varepsilon>0$ there exists $N$ such that
$\|\psi_n-\psi_m\|<\varepsilon$ for all $n,m\geq N$. The space $V$
is \emph{complete} if every Cauchy sequence in $V$ converges to a
limit in $V$. A complete inner-product space is, by definition, a
\emph{Hilbert space}.
\end{definition}

\begin{definition}[Separability and orthonormal basis]
\label{def:separable}
A topological space is \emph{separable} if it contains a countable
dense subset. For a Hilbert space, separability is equivalent to
the existence of a countable \emph{orthonormal basis}: a sequence
$\{\ket{n}\}_{n\in\N}$ with $\braket{n}{m}=\delta_{nm}$ such that
every $\ket{\psi}\in\hilbert$ admits the (norm-convergent) expansion
$\ket{\psi}=\sum_n c_n\ket{n}$ with $c_n=\braket{n}{\psi}$. Every
Hilbert space encountered in the main text is separable.
\end{definition}

% --------------------------------------------------------------
\section{Dual space and operators}
\label{app:dual-space}

\begin{definition}[(Continuous) dual space]
\label{def:dual-space}
The \emph{dual space} $V^*$ of a Hilbert space $V$ is the space of
continuous linear functionals $\phi:V\to\C$. The Riesz
representation theorem identifies $V^*$ with $V$ itself: for every
$\phi\in V^*$ there is a unique $\chi_\phi\in V$ with
$\phi(\psi)=\langle\chi_\phi,\psi\rangle$. With the physics
convention used in this document (\cref{def:sesquilinear},
antilinear in the first slot), the resulting map
$\chi\mapsto\bra{\chi}$ is itself \emph{anti-linear} in $\chi$:
$\bra{\alpha\chi+\beta\eta}=\bar\alpha\bra{\chi}+\bar\beta\bra{\eta}$.
Dirac notation makes the identification explicit:
$\bra{\phi}\in V^*$ is the linear functional that sends
$\ket{\psi}$ to $\braket{\phi}{\psi}$.
\end{definition}

\begin{definition}[Linear operator, bounded operator]
\label{def:linear-op}
A map $\op{A}:V\to V$ is a \emph{linear operator} if
$\op{A}(\alpha\phi+\beta\psi)=\alpha\op{A}\phi+\beta\op{A}\psi$ for
all $\alpha,\beta\in\C$ and $\phi,\psi\in V$. It is \emph{bounded}
if $\|\op{A}\|\equiv\sup_{\|\psi\|=1}\|\op{A}\psi\|<\infty$. In
finite dimensions every linear operator is bounded; in infinite
dimensions important operators (e.g.\ $\op{x}$ and $\op{p}$) are
unbounded and require domain restrictions to be made rigorous.
\end{definition}

\begin{definition}[Adjoint]
\label{def:adjoint}
The \emph{adjoint} (or Hermitian conjugate) $\op{A}^\dagger$ of a
bounded linear operator $\op{A}$ on a Hilbert space is the unique
bounded operator satisfying
\begin{equation}
\langle\op{A}^\dagger\phi,\psi\rangle = \langle\phi,\op{A}\psi\rangle
\qquad\text{for all}\ \phi,\psi\in V.
\end{equation}
The Hermitian, unitary, and projector classes of operators are
defined in \cref{def:operator-types} in terms of the relation
between $\op{A}$ and $\op{A}^\dagger$.
\end{definition}

% --------------------------------------------------------------
\section{Tensor products}
\label{app:tensor}

\begin{definition}[Tensor product of Hilbert spaces]
\label{def:tensor-product}
The \emph{tensor product} $V\otimes W$ of two Hilbert spaces is
the (completion of the) vector space of finite formal sums
$\sum_i \alpha_i\,\phi_i\otimes\chi_i$ with $\phi_i\in V$,
$\chi_i\in W$, $\alpha_i\in\C$, modulo bilinearity in each
factor. Its inner product, defined on product elements by
\begin{equation}
\langle\phi\otimes\chi,\,\phi'\otimes\chi'\rangle
= \langle\phi,\phi'\rangle_V\,\langle\chi,\chi'\rangle_W,
\end{equation}
extends sesquilinearly to all of $V\otimes W$. If
$\{\ket{i}\}_i$ and $\{\ket{j}\}_j$ are orthonormal bases of $V$
and $W$ then $\{\ket{i}\otimes\ket{j}\}_{ij}$ is an orthonormal
basis of $V\otimes W$.
\end{definition}
